% from ega0-15-fr.tex
\oldpage[0\textsubscript{IV-1}]{15}
dans $M'_k/M'_{k+1}$. On est donc ramené à écrire que, pour $x\in M_i$, la
relation
\[
  f^{k-i}x\in M_k+fM_{k-1}+\cdots+
  f^{k-i-1}M_{i+1}+M'_{k+1}
  \tag{15.1.6.2}\label{0.15.1.6.2-fr}
\]
implique $x\in fM_i+M_{i+1}$. Or, le second membre de (15.1.6.2) est contenu
dans $M_{i+1}+M'_{k+1}$, et en vertu de (15.1.5.1), $M'_{k+1}$ est contenu
dans $M_{i+1}+f^{k+1-i}M$. L'hypothèse sur $f$ entraîne que $f$ est non
diviseur de zéro dans $M/M_h$ pour tout $h$ (Bourbaki, \emph{loc.~cit.}). La
relation $f^{k-i}x\in f^{k+1-i}M+M_{i+1}$ entraîne donc (en raisonnant dans
$M/M_{i+1}$) l'existence d'un $y\in M$ tel que $x-fy\in M_{i+1}$, et comme
$x\in M_i$, on a $fy\in M_i$, d'où (en raisonnant dans $M/M_i$), $y\in M_i$,
et finalement $x\in fM_i+M_{i+1}$. C.Q.F.D.
\end{proof}

\begin{definition}[15.1.7]
\label{0.15.1.7-fr}
Soit $(f_i)_{1\leq i\leq n}$ une suite d'éléments de $A$. On dit qu'elle est
\emph{$M$-régulière} si, pour $1\leq i\leq n$, $f_i$ est
$\left(M/(\sum_{1\leq j\leq i-1}f_jM)\right)$-régulier. On dit que la suite
$(f_i)$ est \emph{$M$-quasi-régulière} si
l'homomorphisme canonique $\varphi$ défini dans (15.1.1.1) est bijectif.

Lorsque $M=A$, on dit simplement «~suite régulière~» et
«~suite quasi-régulière~».
\end{definition}

\begin{proposition}[15.1.8]
\label{0.15.1.8-fr}
Les notations étant celles de (15.1.5), on suppose que la suite
$(f_i)_{1\leq i\leq n}$ est $\gr_\bullet(M)$-régulière. Alors
l'homomorphisme canonique (15.1.5.2) est bijectif.
\end{proposition}

\begin{proof}
La proposition n'est autre que (15.1.6) pour $n=1$~; raisonnons par récurrence
sur $n$. Posons
\[
  \mathfrak{J}''=\sum_{i=1}^{n-1}f_iA,
\]
de sorte que $\mathfrak{J}=\mathfrak{J}''+f_nA$, et soit
\[
  M''_k=M_k+\mathfrak{J}''M_{k-1}+\cdots+
  \mathfrak{J}''^{k-1}M_1+\mathfrak{J}''^kM_0~;
\]
on peut donc écrire
\[
  M'_k=M''_k+f_nM''_{k-1}+\cdots+f_n^{k-1}M''_1+f_n^kM''_0.
\]
Notons $\gr''_\bullet(M)$ le $A$-module gradué associé à la filtration
$(M''_k)$ de $M$. On peut écrire, à des isomorphismes canoniques près
\[
  (\gr_\bullet(M)\otimes(A/\mathfrak{J}))\otimes_A A[T_1,\ldots,T_n]
  =
  (\gr_\bullet(M)\otimes(A/\mathfrak{J}''))[T_1,\ldots,T_{n-1}]
  \otimes(A/f_nA)\otimes_A A[T_n]
\]
et par suite (15.1.5.2) se factorise en
\[
  (\gr''_\bullet(M)\otimes(A/f_nA))[T_n]\longrightarrow\gr'_\bullet(M)
  \tag{15.1.8.1}\label{0.15.1.8.1-fr}
\]
et
\[
  ((\gr_\bullet(M)\otimes(A/\mathfrak{J}''))[T_1,\ldots,T_{n-1}])
  \otimes(A/f_nA)\otimes_A A[T_n]
  \xrightarrow{\psi'\otimes1}
  \gr''_\bullet(M)\otimes(A/f_nA)\otimes_A A[T_n]
\]
$\psi'$ étant l'application (15.1.5.2) où l'on remplace $n$ par $n-1$,
$\mathfrak{J}$ par $\mathfrak{J}''$ et $\gr'_\bullet(M)$ par
$\gr''_\bullet(M)$. L'hypothèse de récurrence entraîne que $\psi'$ est
bijective et il reste donc à voir qu'il en est de même de (15.1.8.1). En
appliquant (15.1.6), il suffit de montrer que $f_n$ est
$\gr''_\bullet(M)$-régulier. Mais $\gr''_\bullet(M)$ s'identifie à
\[
  (\gr_\bullet(M)\otimes(A/\mathfrak{J}''))
  [T_1,\ldots,T_{n-1}]~;
\]
par hypothèse, l'homothétie de rapport $f_n$ dans
$\gr_\bullet(M)\otimes(A/\mathfrak{J}'')$ est bijective~; il en est par suite
de même de l'homothétie de rapport $f_n$ dans $\gr''_\bullet(M)$, puisque
$A[T_1,\ldots,T_{n-1}]$ est un $A$-module libre.
\end{proof}

\begin{proposition}[15.1.9]
\label{0.15.1.9-fr}
Soient $M$ un $A$-module, $(f_i)_{1\leq i\leq n}$ une suite d'éléments de $A$.
Si la suite $(f_i)$ est $M$-régulière, elle est $M$-quasi-régulière. La
réciproque est vraie lorsqu'on suppose en outre que les $A$-modules
\[
  M,\quad M/f_1M,\quad\ldots,\quad
  M/\left(\sum_{1\leq j\leq n-1}f_jM\right)
\]
sont séparés pour la topologie $\mathfrak{J}$-préadique (où
$\mathfrak{J}=\sum_{i=1}^n f_iA$).
\end{proposition}

\begin{proof}
La première assertion résulte de (15.1.8), où l'on prend $M_1=0$, et par suite
l'homomorphisme (15.1.5.2) se réduit à (15.1.1.1). Inversement, supposons que
l'homomorphisme $\varphi$ de (15.1.1.1) soit bijectif. Notons que si on pose
\[
  \mathfrak{J}''=\sum_{i=1}^{n-1}f_iA,
\]
